\documentclass[10pt,a4paper]{article}
\usepackage[margin=22mm]{geometry}
\usepackage{amsmath,amssymb,booktabs,array}
\usepackage[T1]{fontenc}
\usepackage{lmodern}
\usepackage[hidelinks]{hyperref}
\setlength{\parindent}{0pt}
\setlength{\parskip}{4pt}
\title{EE3032 Introduction to VLSI Design\\2024 Homework 4: Pre-layout SPICE Deliverable}
\author{Independent reconstruction and verification}
\date{13 October 2026}
\begin{document}
\maketitle

\section{Requested scope and relation to the original assignment}
The original Homework 4 asks for layout, extracted netlists and a pre-/post-layout
comparison. This deliverable follows the user's explicit request to stop at the
\texttt{.sp} stage. It provides a transistor-level schematic netlist and a runnable
testbench. No layout, DRC, LVS, PEX or post-layout simulation is claimed.

The final Homework 3 netlist is copied without modifying transistor connections
or dimensions. The copied file is \texttt{four\_bit\_multbased\_adder.sp}; its hash
and the new testbench hash are recorded in \texttt{statistics.json}. This makes
the correspondence between the verified design and the Homework 4 deliverable
explicit.

\section{Function and interface}
For bit $i$, define the equality signal $e_i=\overline{x_i\oplus y_i}$. The cell
implements the two XNOR blocks and carry multiplexer in the assignment figure:
\[
 s_i=\overline{e_i\oplus c_i}=x_i\oplus y_i\oplus c_i,\qquad
 c_{i+1}=e_i y_i+\overline{e_i}c_i.
\]
If $x_i=y_i$, the carry is $y_i$; if $x_i\ne y_i$, the carry is $c_i$.
Four cells implement
\[
 \{\mathrm{COUT},S_3,S_2,S_1,S_0\}=X+Y+\mathrm{CIN}.
\]
The public subcircuit pin order is
\begin{center}\small
\texttt{CIN X0 X1 X2 X3 Y0 Y1 Y2 Y3 S0 S1 S2 S3 COUT}.
\end{center}
The least significant bits are \texttt{X0}, \texttt{Y0}, \texttt{S0}.
The testbench defines global rails \texttt{vdd} and \texttt{vss}.

\section{Simulation and numerical checks}
The external process file is the user's \texttt{D:/cic018.l}, TT corner,
$V_{DD}=1.8$ V and $T=25^\circ$C. The model is not included in the repository.
The testbench preserves the instructor's PWL sources and CIN=0, simulating
0--9 ns with no extra output load. The sampling points are 1.8, 2.8, $\ldots$,
8.8 ns, after each input vector has completed its transition.

\begin{center}
\begin{tabular}{rrrr}
\toprule
Sample (ns)&$X$&$Y$&Expected $\{C_{out},S\}$ (decimal)\\
\midrule
1.8&1&5&6\\2.8&15&13&28\\3.8&11&9&20\\4.8&4&8&12\\
5.8&3&7&10\\6.8&12&14&26\\7.8&2&15&17\\8.8&6&10&16\\
\bottomrule
\end{tabular}
\end{center}

The runnable verification records all five output voltages for all eight vectors.
For an expected zero it checks $-0.02\le V\le0.2$ V; for an expected one it checks
$1.6\le V\le1.82$ V. These tolerance bands allow small numerical overshoot while
requiring output levels near both rails. The complete voltages are in
\texttt{vectors.csv}; this numerical test does not use a GUI screenshot.

HSPICE's \texttt{power} includes dissipation and changes in stored energy. It is
the negative sum of the powers absorbed by all independent sources, including
the PWL input drivers. VDD-only power is measured separately as
$-V(\mathrm{vdd})I(\mathrm{vvdd})$. The measured total energy must satisfy
\[
 E=\int_0^{9\,\mathrm{ns}}P(t)\,dt
   =P_{\mathrm{avg}}\,(9\,\mathrm{ns}).
\]
The file \texttt{statistics.json} records power, energy, pass counts,
simulator version and source hashes. Homework 3 documents delay and exhaustive
input checks for the same netlist. These are pre-layout results only.

\begin{center}
\begin{tabular}{lr}
\toprule
Final netlist, independently rerun for this folder & Result\\
\midrule
Transistor count & 64\\
All widths / all lengths & $0.47/0.18\,\mu$m\\
$\sum_i W_iL_i$ (schematic WL metric) & $5.4144\,\mu\mathrm{m}^2$\\
Total average HSPICE power & $69.558498\,\mu$W\\
Supply-delivered average power & $37.4328403\,\mu$W\\
Total energy over 9 ns & $626.026482$ fJ\\
Assigned-vector checks & 8 / 8 pass\\
\bottomrule
\end{tabular}
\end{center}
The power and energy values above use a 1 ps maximum timestep. Homework 3
separately records the original 0.1 ns-output-step grading testbench and its
slightly different adaptive-step numerical results. Its full 512-case
transient check applies to the identical transistor netlist supplied here.

\section{Review of the existing Homework 4 answer}
The supplied previous answer contains a 64-transistor schematic and extracted
parasitic files. Its schematic was rerun as a separate DUT using the exact
instructor stimulus. It also passed all 512 DC combinations of
$(X,Y,\mathrm{CIN})$ using the 0.2/1.6 V output criteria. This is evidence that
the reference is functional under the tested conditions. The refined energy,
delay and area comparison, using a 1 ps maximum timestep, is given in Homework 3.
No unmeasured corner or load robustness is inferred.

The previous report reverses the pre-/post-layout labels: the extracted circuit
with parasitics is the post-layout circuit. Its testbench simultaneously
instantiates the schematic and the extracted design on a shared supply;
\texttt{avg power} therefore combines both DUTs. The report's energy and
layout area cannot be reused as measurements of the current design. The current
testbench instantiates exactly one schematic DUT. The design comparison and
selection are recorded in the Homework 3 report and statistics tables.

\section{Reproduction and area distinction}
Run \texttt{python reproduce.py} in this folder, or set \texttt{--hspice} and
\texttt{--model} to the local simulator and process paths. Simulator raw outputs
remain in \texttt{run\_artifacts/}; the small statistics files are published.
Licensed access to HSPICE and access to the external process file are required.

Homework 3 uses $A_{\mathrm{WL}}=\sum_i W_iL_i$. Homework 4's physical layout
area is the rectangle covering all layout materials. These measures are not
interchangeable. Since this scope includes no physical layout, no physical
area or post-layout delay/power is reported.
\end{document}
